% !TeX root = ../main.tex
% !TeX encoding = utf8
	\section{CƠ SỞ TOÁN HỌC CỦA ECDF VÀ THUẬT TOÁN MONTE CARLO}
	
	\subsection{Phân phối thực nghiệm (\textit{Empirical Distribution}) và các đặc trưng đại số}
	Cho $x = (x_1, x_2, \dots, x_n)$ là một mẫu ngẫu nhiên gồm $n$ quan sát độc lập cùng phân phối (i.i.d.) rút từ hàm phân phối tích lũy $F(t) = \mathbb{P}(X \le t)$ chưa biết của tổng thể \cite{buithuytrang2026}.
	
	\begin{definition}[\textbf{Hàm phân phối thực nghiệm ECDF}]
		Hàm phân phối thực nghiệm $F_n^*(t)$ (hoặc $F_x^*(t)$) của mẫu quan sát $x$ được định nghĩa là hàm nấc thang đo tỷ lệ các quan sát trong mẫu nhỏ hơn hoặc bằng $t$ \cite{buithuytrang2026}:
		\begin{equation}
			F_n^*(t) = \frac{1}{n}\sum_{i=1}^{n} \mathbb{I}_{(-\infty, t]}(x_i) = \frac{\#\{x_i \le t\}}{n}, \quad \forall t \in \mathbb{R}
		\end{equation}
		trong đó $\mathbb{I}_A(\cdot)$ là hàm chỉ thị (indicator function):
		\begin{equation}
			\mathbb{I}_A(y) = \begin{cases} 
				1 & \text{nếu } y \in A \\ 
				0 & \text{nếu } y \notin A 
			\end{cases}
		\end{equation}
	\end{definition}
		Xét mẫu minh họa $x=\{12,15,18,21,24\}$. Tại mỗi giá trị quan sát
		$x_{(i)}$, hàm ECDF bằng tỷ lệ số quan sát không vượt quá giá trị đó:
		
		\begin{table}[htbp]
			\centering
			\caption{Tính hàm phân phối thực nghiệm từ một mẫu}
			\label{tab:ecdf-values}
			\begin{tabular}{ccc}
				\toprule
				Giá trị $x_{(i)}$ & Số quan sát $\leq x_{(i)}$ & $F_n^*(x_{(i)})$ \\
				\midrule
				12 & 1 & 0.2 \\
				15 & 2 & 0.4 \\
				18 & 3 & 0.6 \\
				21 & 4 & 0.8 \\
				24 & 5 & 1.0 \\
				\bottomrule
			\end{tabular}
		\end{table}
	
	\textbf{Các tính chất xác suất của $F_n^*(t)$:}
	\begin{enumerate}
		\item \textbf{Tính chất phân phối nhị thức:}
		Trước khi quan sát dữ liệu, ký hiệu $X_1,\dots,X_n$ là các biến ngẫu
		nhiên độc lập cùng phân phối theo $F$. Với một giá trị $t$ cố định, đặt
		$Y_i=\mathbf{1}\{X_i\leq t\}$. Khi đó $Y_i$ là biến Bernoulli với xác
		suất thành công $F(t)$, nên:
		\begin{equation}
			nF_n^*(t)=\sum_{i=1}^{n}Y_i
			\sim \operatorname{Binomial}\bigl(n,F(t)\bigr).
		\end{equation}
		\item \textbf{Tính không chệch (Unbiasedness):} Kỳ vọng của $F_n^*(t)$ đúng bằng phân phối tổng thể thực sự $F(t)$ tại mọi điểm $t$ \cite{buithuytrang2026}:
		\begin{equation}
			\mathbb{E}[F_n^*(t)] = \mathbb{E}\left[\frac{1}{n}\sum_{i=1}^n Y_i\right] = \frac{1}{n} \cdot n F(t) = F(t)
		\end{equation}
		\item \textbf{Phương sai và hội tụ tại từng điểm:} Phương sai của $F_n^*(t)$ tiến về 0 khi $n \to \infty$ \cite{buithuytrang2026}:
		\begin{equation}
			\text{Var}\left(F_n^*(t)\right) = \text{Var}\left(\frac{1}{n}\sum_{i=1}^n Y_i\right) = \frac{n F(t)(1 - F(t))}{n^2} = \frac{F(t)(1 - F(t))}{n} \xrightarrow{n \to \infty} 0
		\end{equation}
		Với mỗi $t$ cố định, bất đẳng thức Chebyshev cho thấy:
		\begin{equation}
			\mathbb{P}\left(\left|F_n^*(t)-F(t)\right|>\varepsilon\right)
			\leq
			\frac{F(t)(1-F(t))}{n\varepsilon^2}
			\longrightarrow 0.
		\end{equation}
		Vì vậy, với mỗi $t$ cố định, $F_n^*(t)$ hội tụ theo xác suất về $F(t)$.
		\item \textbf{Định lý Glivenko-Cantelli:} Hàm phân phối thực nghiệm $F_n^*(t)$ hội tụ đều về hàm phân phối tổng thể $F(t)$ hầu chắc chắn trên toàn bộ đường thẳng thực $\mathbb{R}$ \cite{efron1993}:
		\begin{equation}
			\mathbb{P}\left( \lim_{n \to \infty} \sup_{t \in \mathbb{R}} |F_n^*(t) - F(t)| = 0 \right) = 1
		\end{equation}

	\newpage
		\item \textbf{Minh họa bằng mô phỏng R:} Mô phỏng dưới đây lấy mẫu từ phân phối Chuẩn $N(0,1)$ với các cỡ mẫu $n=10,30,100,1000$, sau đó so sánh ECDF với hàm phân phối lý thuyết $F(t)=\Phi(t)$.

	\begin{lstlisting}[style=Rstyle, caption={Mô phỏng ECDF}]
		# 1. Thiet lap tham so va phan phoi tong the N(0, 1)
		set.seed(2026)
		n_sizes <- c(10, 30, 100, 1000)
		colors  <- c("red", "orange", "blue", "green3")
		
		# 2. Ve CDF ly thiet
		curve(pnorm(x), from = -3, to = 3, col = "black", lwd = 3, lty = 1,
		main = "Minh hoa Dinh ly Glivenko-Cantelli (Hoi tu ECDF -> CDF)",
		xlab = "t", ylab = "F(t)")
		
		# 3. Ve ECDF cho cac kich thuoc mau n khac nhau
		for (i in 1:length(n_sizes)) {
			n <- n_sizes[i]
			sample_data <- rnorm(n)
			fn <- ecdf(sample_data)
			plot(fn, add = TRUE, col = colors[i], lwd = 2, do.points = FALSE)
		}
		
		legend("topleft", 
		legend = c("CDF Ly thiet N(0,1)", "n = 10", "n = 30", "n = 100", "n = 1000"),
		col = c("black", colors), lty = 1, lwd = 2)
	\end{lstlisting}
	\end{enumerate}
	
	
	\subsection{Nguyên lý Plug-in và Mô hình Không gian Xác suất Bootstrap}
	
	\begin{definition}[\textbf{Nguyên lý Plug-in}]
		Giả sử tham số tổng thể cần nghiên cứu có thể biểu diễn dưới dạng một phiếm hàm (functional) $\theta = \theta(F)$ phụ thuộc vào hàm phân phối tổng thể $F$ \cite{buithuytrang2026}. Ước lượng Plug-in $\hat{\theta}_n$ của $\theta$ được xác định bằng cách thay thế hàm phân phối $F$ chưa biết bằng hàm phân phối thực nghiệm $F_n^*$ \cite{buithuytrang2026}:
		\begin{equation}
			\hat{\theta}_n = \theta(F_n^*)
		\end{equation}
	\end{definition}
	
	\textbf{Khai triển Đại số cho các Ước lượng Plug-in Thường gặp:}
	\begin{itemize}
		\item \textbf{Trung bình tổng thể $\mu(F) = \int x dF(x)$:} Thay $F$ bằng $F_n^*$, ta thu được trung bình mẫu $\bar{x}$ \cite{buithuytrang2026}:
		\begin{equation}
			\hat{\mu}_{\text{plug-in}} = \mu(F_n^*) = \int x dF_n^*(x) = \sum_{i=1}^n x_i \cdot \mathbb{P}_X^*(X^* = x_i) = \frac{1}{n}\sum_{i=1}^n x_i = \bar{x}
		\end{equation}
		\item \textbf{Phương sai tổng thể $\sigma^2(F) = \int (x - \mu)^2 dF(x)$:} Thay $F$ bằng $F_n^*$ và $\mu$ bằng $\bar{x}$, ta thu được phương sai thực nghiệm mẫu (mẫu số $n$) \cite{buithuytrang2026}:
		\begin{equation}
			\hat{\sigma}^2_{\text{plug-in}} = \sigma^2(F_n^*) = \int (x - \bar{x})^2 dF_n^*(x) = \frac{1}{n}\sum_{i=1}^n (x_i - \bar{x})^2
		\end{equation}
	\end{itemize}
	
	\newpage

	\textbf{Sơ đồ Tương đương trong "Thế giới Bootstrap":}
	Trong cơ chế Bootstrap, phân phối thực nghiệm $F_n^*$ đóng vai trò phân phối tổng thể thay thế; mẫu quan sát $x$ vì vậy được xem như một tổng thể giả lập \cite{efron1993}:
	\begin{equation}
		\begin{array}{r c l c l c l}
			\textbf{Thế giới thực:}
			& F
			& \xrightarrow{\text{lấy mẫu ngẫu nhiên}}
			& x
			& \xrightarrow{\text{tính thống kê}}
			& \hat{\theta}_n
			& \\[1.2em]
			
			\textbf{Thế giới Bootstrap:}
			& F_n^*
			& \xrightarrow{\text{lấy mẫu có hoàn lại}}
			& x^*
			& \xrightarrow{\text{tính thống kê}}
			& \hat{\theta}_n^*
			&
		\end{array}
	\end{equation}
	
	\subsection{Mẫu Bootstrap và Thuật toán Lấy mẫu lại Monte Carlo}
	
	Một \textbf{Mẫu Bootstrap} $x^* = (x_1^*, x_2^*, \dots, x_n^*)$ là một mẫu ngẫu nhiên kích thước $n$ được rút ngẫu nhiên có hoàn lại từ mẫu quan sát gốc $x = (x_1, \dots, x_n)$ \cite{efron1993}. Mọi phần tử $x_i^*$ đều nhận giá trị $x_j$ bất kỳ trong mẫu gốc với xác suất đồng đều $\mathbb{P}(x_i^* = x_j) = 1/n$ \cite{buithuytrang2026}. Giá trị $\hat{\theta}_n^* = t(x^*)$ tính trên mẫu Bootstrap được gọi là một \textbf{Bản sao Bootstrap (Bootstrap Replicate)}.
	
	Nếu $n$ quan sát ban đầu đôi một khác nhau, số mẫu Bootstrap khác
	nhau khi xem mỗi mẫu như một đa tập, không xét thứ tự, là
	$\binom{2n-1}{n}$. Với $n=20$, con số này bằng
	$\binom{39}{20}=68{,}923{,}264{,}410$. Số lượng lớn này cho thấy việc
	liệt kê toàn bộ các cấu hình thường không thuận tiện, vì vậy ta dùng
	mô phỏng Monte Carlo để xấp xỉ phân phối Bootstrap \cite{efron1993}.
	
	\begin{algorithm}[H]
	\Input{Mẫu dữ liệu gốc $x=(x_1,\dots,x_n)$; số lần lấy mẫu lại $B$; hàm thống kê $t(\cdot)$ \cite{buithuytrang2026}.}
	\Output{$B$ bản sao Bootstrap $\{\hat{\theta}^{*(1)},\dots,\hat{\theta}^{*(B)}\}$ và sai số chuẩn Bootstrap $\widehat{se}^*$ \cite{buithuytrang2026}.}
		\BlankLine
		Tính giá trị ước lượng điểm trên mẫu gốc: $\hat{\theta}_n = t(x)$ \cite{buithuytrang2026}\;
		\For{$b = 1, 2, \dots, B$}{
			Rút ngẫu nhiên có hoàn lại $n$ phần tử từ mẫu $x$ để tạo mẫu Bootstrap $x^{*(b)} = (x_1^{*(b)}, \dots, x_n^{*(b)})$ \cite{buithuytrang2026}\;
			Tính giá trị bản sao Bootstrap thứ $b$: $\hat{\theta}^{*(b)} = t\left(x^{*(b)}\right)$ \cite{buithuytrang2026}\;
		}
		Tính trung bình đại số của $B$ bản sao Bootstrap: $\bar{\theta}^* = \frac{1}{B}\sum_{b=1}^{B} \hat{\theta}^{*(b)}$ \cite{buithuytrang2026}\;
		Tính Ước lượng Sai số chuẩn Bootstrap $\widehat{se}^*$:
		\begin{equation}
			\widehat{se}^* = \sqrt{\frac{1}{B-1}\sum_{b=1}^{B} \left(\hat{\theta}^{*(b)} - \bar{\theta}^*\right)^2} \quad \cite{buithuytrang2026}
		\end{equation}\
		\Return Tập hợp các bản sao $\{\hat{\theta}^{*(b)}\}_{b=1}^B$ và $\widehat{se}^*$ \cite{buithuytrang2026}\;
		\caption{Thuật toán lấy mẫu lại Monte Carlo cho Bootstrap}
	\end{algorithm}
	\subsubsection{Ảnh hưởng của số lần lấy mẫu Bootstrap $B$}
	
	Sau khi cố định mẫu gốc, ký hiệu phân phối có điều kiện của bản sao
	Bootstrap là
	\[
	G_n^*(t)
	=
	\mathbb{P}^*\left(\hat{\theta}_n^* \leq t \mid x\right).
	\]
	Trong tính toán, ta chỉ tạo $B$ bản sao Bootstrap nên ước lượng phân phối
	này là
	\[
	\widehat{G}_{n,B}^*(t)
	=
	\frac{1}{B}\sum_{b=1}^{B}
	\mathbf{1}\left\{\hat{\theta}^{*(b)}\leq t\right\}.
	\]
	
	Với mẫu gốc đã cố định, các chỉ báo trong tổng có phân phối Bernoulli
	với xác suất thành công $G_n^*(t)$. Do đó:
	\begin{align}
		\mathbb{E}^*\left[\widehat{G}_{n,B}^*(t)\right]
		&=G_n^*(t),\\
		\operatorname{Var}^*\left[\widehat{G}_{n,B}^*(t)\right]
		&=\frac{G_n^*(t)\left[1-G_n^*(t)\right]}{B}.
	\end{align}
	
	Khi tăng $B$, phương sai Monte Carlo của $\widehat{G}_{n,B}^*(t)$ giảm.
	Điều này làm xấp xỉ phân phối Bootstrap có điều kiện ổn định hơn, nhưng
	không khắc phục được việc mẫu gốc thiếu tính đại diện hoặc cách lấy mẫu
	lại không phù hợp với cấu trúc dữ liệu. Sai số mô phỏng do $B$ hữu hạn
	cần được phân biệt với sai số của phép xấp xỉ Bootstrap.
	
\subsubsection{Điều kiện để phép xấp xỉ Bootstrap có ý nghĩa}

Bootstrap thay phân phối tổng thể chưa biết $F$ bằng phân phối thực
nghiệm $F_n^*$. Tuy nhiên, việc thay thế này không bảo đảm mọi phân phối
Bootstrap đều xấp xỉ tốt phân phối lấy mẫu của mọi thống kê. Ta cần xét
điều kiện áp dụng và tính chất của thống kê đang nghiên cứu
\cite{bickelfreedman1981}.

Xét trường hợp thống kê là trung bình mẫu. Giả sử
$X_1,\ldots,X_n$ độc lập cùng phân phối, có kỳ vọng $\mu$ và phương sai
hữu hạn, dương $\sigma^2$. Ký hiệu $\bar{X}^*$ là trung bình của một mẫu
Bootstrap và $s_n$ là độ lệch chuẩn mẫu. Khi đó, dưới các điều kiện trên,
phân phối có điều kiện của thống kê Bootstrap chuẩn hóa hội tụ về phân
phối Chuẩn tắc:
\begin{equation}
	\sup_{z\in\mathbb{R}}
	\left|
	\mathbb{P}^*\left(
	\frac{\sqrt{n}(\bar{X}^*-\bar{X})}{s_n}\leq z
	\right)
	-\Phi(z)
	\right|
	\xrightarrow{P} 0,
\end{equation}
trong đó $\mathbb{P}^*$ là xác suất có điều kiện khi giữ cố định mẫu quan
sát ban đầu, còn sự hội tụ theo xác suất được xét theo cơ chế lấy mẫu ban
đầu. Kết quả này là cơ sở để dùng các bản sao Bootstrap xấp xỉ phân phối
lấy mẫu của trung bình mẫu trong các điều kiện đã nêu.

Kết quả trên chỉ áp dụng cho trung bình mẫu dưới các điều kiện đã nêu.
Với mẫu nhỏ, phân phối thực nghiệm có thể chưa thể hiện tốt phần đuôi của
tổng thể; với dữ liệu phụ thuộc, cách lấy mẫu lại cần giữ cấu trúc phụ
thuộc. Khi áp dụng Bootstrap, cần xác định thống kê quan tâm, giả định về
dữ liệu và đơn vị lấy mẫu lại
\cite{bickelfreedman1981, davison1997}.
	
	% ==================== PHẦN 3 ====================
